\section{Introduction}
\subsection{Overview}
An experiment may continue while new actions become available. A product comparison can admit new designs; an allocation system can receive new candidates from an external development process. Existing observations remain informative about old actions, but each arrival introduces an alternative for which the system has no data. A natural objective is to limit the cumulative loss relative to the best action available at each decision time. We study this objective when actions remain available after entry and their reward distributions do not change.

The final number of actions gives an incomplete description of this problem. Consider an incumbent and $M$ new actions over $T$ rounds. If all new actions are present from the beginning, the usual worst-case learning cost is of order $\min\{T,\sqrt{MT}\}$. If the same actions arrive in the last round, the regret is at most one. With several arrival batches, neither the longest remaining lifetime nor the sum of the lifetimes generally determines the rate. Learning opportunities overlap, and several time scales can each contribute appreciably to the total loss. The question is therefore which feature of the entire arrival schedule measures its statistical difficulty.

We answer this question for every deterministic exogenous schedule. Each arm has an independent, stationary reward distribution supported on $[0,1]$. Arrivals can occur in arbitrarily large finite batches, and ties are allowed. We fix the schedule before maximizing expected regret over reward distributions. The benchmark chooses the available arm with largest mean at each round. In particular, the benchmark changes only because of arrivals, not because an existing reward distribution drifts.

The answer is a sum over remaining-lifetime scales. For a fixed endpoint, group the competing arms whose remaining lifetimes are between $2^{\ell-1}$ and $2^\ell$. A group containing $k_\ell$ arms contributes $\min\{2^\ell,\sqrt{k_\ell2^\ell}\}$. We show that the sum of these contributions equals the schedule-specific minimax regret up to universal multiplicative constants. A learner that knows neither the endpoint nor future arrivals achieves this rate. It uses a MOSS index with a clock that restarts when the number of competing arms doubles, while its sample counts and reward sums continue across every restart.

The main difficulty in the lower bound is to make the scale-specific testing problems coexist. Adding separate worst-case risks is invalid if their respective hard environments cannot occur together. Our construction selects a constant fraction of the relevant scale mass and places independent testing problems in disjoint time windows. The benchmark increments are summable, and the reference means do not reveal the outcomes of earlier tests. Information acquired before a testing window is included in its likelihood calculation. The upper bound faces a different issue: retaining observations prevents a conditional restart argument. We instead control all potential sample prefixes and charge each arm-count pair at most once.

A second expression makes the complexity computable without a dyadic grid. We maximize a clipped square-root objective over the fractional sample allocations allowed by the arrivals. The constraints are nested suffix capacities, and their optimizer has a lower-convex-hull representation. This is an instance of established nested resource allocation; its statistical role here is to identify a quantity equivalent to regret. The representation also gives exact stability properties, a sharp comparison with integer sample allocations and a precise explanation of diminishing marginal capacity when more arms are added.

The results distinguish several notions that are otherwise easy to conflate. A natural lifetime fixed point measures the Euclidean norm of the scale costs, while the minimax risk measures their sum. The two can differ by a factor of order $\sqrt{\log T}$. Retaining observations can reduce regret from squared-logarithmic to logarithmic order relative to an otherwise identical policy that erases them, even though both policies have the same minimax rate.

Random arrivals introduce another distinction. The reward environment must be fixed before the schedule is drawn, so separate least-favorable environments for different schedules cannot be combined. We first exhibit a law under which the average deterministic complexity grows logarithmically while joint minimax risk stays bounded. We then characterize the joint risk for fixed batch templates whose relative times share one random stretch. A truncated-lifetime functional identifies the amount of difficulty that a common environment can sustain across stretches. The endpoint is known in this result, and the first batch reveals the remaining schedule. The constants depend on the fixed template; the result does not address growing templates or arbitrary arrival laws.

The numerical studies follow these distinctions. Compatible product priors connect the multiscale lower construction to observed regret, with an exact Bayes calculation at small horizons. A factorial comparison separates clock changes from capacity calibration, while a matched pair isolates the effect of forgetting. Independent development and evaluation samples compare one reward environment across random schedules with environments chosen separately for each schedule. The studies include failures of the proposed index and retain the information differences between the policies being compared.

\subsection{Relation to Existing Work}
Stochastic sleeping bandits already allow changing availability sets. Ordering arms by their means makes the best-ordering comparator coincide with the best currently available arm. The AUER analysis of \citet{kleinberg2008sleeping} gives logarithmic bounds involving adjacent mean gaps for a finite pool and retains observations between appearances. Thus neither the comparator nor data retention is specific to the present model. \citet{gaillard2023sleeping} further relate sleeping internal regret to ordering regret under independent and identically distributed losses. Per-action guarantees such as those of \citet{nguyen2024sleeping} use another comparator: a fixed action on the rounds when it is active. These results do not by themselves give the constant-factor minimax risk for each specified monotone arrival schedule.

The growing-action setting is studied directly by \citet{ghalme2021ballooning}. Their ballooning-bandit formulation uses the best available arm and establishes both impossibility results and sublinear guarantees under assumptions on the arrival of a best arm, with further discussion of arrival rates. Their BL-MOSS procedure uses the horizon and arrival information to select a collection of arms and subsequently stops admitting new ones. Our characterization conditions on an arbitrary fixed schedule and maximizes over all bounded stationary reward laws. It imposes no tail assumption on the arrival of an optimal arm and requires no cutoff for admitting arms.

The index itself comes from MOSS \citep{audibert2009minimax}. Anytime minimax-optimal MOSS procedures for fixed arm sets were developed by \citet{degenne2016anytime}; the coexistence of minimax and instance-dependent guarantees also has substantial precedent, including \citet{menard2017minimax}. We use this literature as the local learning component. The contribution is the schedule-specific upper and lower characterization, with an upper-bound argument that permits clock recalibration while retaining observations as the arm pool grows. Our gap bound is not an instance-optimality result.

The allocation representation is related to separable convex optimization with nested constraints. \citet{vidal2016nested} explicitly derive lower-convex-hull solutions for shared-shape objectives, and \citet{vidal2019separable} treat nonsmooth convex objectives within a broader class. General submodularity results for concave optimization over polymatroids are available in \citet{he2012polymatroid}. We use these structures with explicit proofs for the present clipped square-root objective. The connection to bandit risk, rather than a new convex optimization algorithm, is the point of the representation.
